\begin{figure}[!ht]
\centering
\begin{tikzpicture}[alt={Actual light-theme native Memory view after the scalar block function. The first byte row starts 10 f1 e7 e4 d1 3b 59 15 and the next row starts c7 d1 f4 c7. An outline identifies the first four bytes.}]
\node[anchor=north west,inner sep=0,font=\small\bfseries] (title) at (0,0) {Words and bytes describe the same output};
\node[anchor=north west,inner sep=0] (a) at ([yshift=-2mm]title.south west) {\MARSRevisionCrop{revision/captures/chacha/chacha-block.png}{0.99\linewidth}{0.03125}{0.31555556}{0.58203125}{0.60888889}};
\begin{scope}[shift={(a.south west)},x={($(a.south east)-(a.south west)$)},y={($(a.north west)-(a.south west)$)}]
\draw[black,line width=.7pt] (.20,.3088) rectangle (.39,.6176);
\end{scope}
\end{tikzpicture}
\caption{The first 32 generated bytes are shown. The first four, 10 f1 e7 e4, encode word e4e7f110 in little-endian order. The outline identifies that word; the full 64-byte output is checked separately.}
\label{fig:revision-chacha-native-block}
\end{figure}
